\documentclass[11pt]{article}
\usepackage[utf8]{inputenc}
\usepackage{amsmath}
\usepackage{amssymb}
\usepackage{geometry}
\usepackage{booktabs}
\usepackage{hyperref}
\geometry{margin=1in}

\hypersetup{colorlinks=true, linkcolor=blue, urlcolor=blue}

\title{MethylValidation Theoretical Foundation}
\author{MethylPipeline}
\date{}

\begin{document}
\maketitle
\tableofcontents
\newpage

% -----------------------------------------------------------------------
\section{Goal}

\texttt{MethylValidation} estimates the \textbf{sampling distribution} of
validation metrics (balanced accuracy, sensitivity, specificity, $F_1$, etc.)
by running the full MethylPipeline many times with different stratified
train\,/\,validation splits (Monte Carlo).

The ``distribution'' MethylValidation provides is the \emph{empirical
distribution} over iterations: with enough iterations, the summary
(mean, std, percentiles) approximates the sampling distribution of each
metric under random train\,/\,validation partitions.

% -----------------------------------------------------------------------
\section{Stratified Train / Validation Split}

For two groups — \emph{control} (e.g.\ healthy) and \emph{disease} —
MethylValidation splits each group independently using the same
\texttt{train\_fraction} $\rho \in (0,1)$:

\begin{align}
  n_{\text{train},k} &= \left\lfloor \rho \cdot n_k \right\rfloor, \\
  n_{\text{val},k}   &= n_k - n_{\text{train},k},
\end{align}

where $k \in \{\text{control, disease}\}$ and $n_k$ is the total sample count
for that group.  The training set is used for the full pipeline
(MethylCentroid $\to$ MethylDetector $\to$ MethylClassifier); the validation
set is held out and used only by \texttt{methyl-predictor}.

\paragraph{Properties.}
\begin{itemize}
  \item Both train and validation sets contain both classes (stratified).
  \item Each iteration uses a different random split (with optional fixed seed
        for reproducibility).
  \item Validation metrics are computed on unseen samples, so they reflect
        generalisation.
\end{itemize}

% -----------------------------------------------------------------------
\section{Per-Iteration Protocol}

Let $\omega = 1, \ldots, n_\mathrm{iter}$ index iterations.

\begin{enumerate}
  \item \textbf{Split.}  Draw a stratified partition from each group.
  \item \textbf{Project.}  Generate a run-specific project JSON with train
        samples only (and optional centroid delta overrides for
        \texttt{add\_samples}/\texttt{remove\_samples}).
  \item \textbf{Pipeline.}  Run
        \texttt{methyl-centroid} $\to$ \texttt{methyl-detector} $\to$
        \texttt{methyl-classifier} on the training data.
  \item \textbf{Validation.}  Run \texttt{methyl-predictor} on the holdout
        validation set; read \texttt{validation\_metrics.json}.
  \item \textbf{Collect.}  Record scalar metrics
        $m_\omega = (\mathrm{BA}_\omega, \mathrm{Sens}_\omega,
                    \mathrm{Spec}_\omega, F_{1,\omega}, \ldots)$
        and step timings.
\end{enumerate}

% -----------------------------------------------------------------------
\section{Empirical Distribution of Metrics}

After $n_\mathrm{iter}$ successful iterations, the collection
$\{m_\omega\}_{\omega=1}^{n_\mathrm{iter}}$ forms an empirical sample of
each metric under random data splits.

\paragraph{Summary statistics.}

For each scalar metric $M$:

\begin{align}
  \bar{M} &= \frac{1}{n_\mathrm{iter}} \sum_\omega M_\omega, \\[4pt]
  s_M &= \sqrt{\frac{1}{n_\mathrm{iter}-1}
               \sum_\omega (M_\omega - \bar{M})^2}, \\[4pt]
  \hat{Q}_p &= \text{empirical } p\text{-th percentile of } \{M_\omega\}.
\end{align}

These are written to \texttt{metrics\_summary.json}.  Key percentiles are
\(p \in \{5, 25, 50, 75, 95\}\).

% -----------------------------------------------------------------------
\section{Processing and Storage Estimation}

\paragraph{Time.}

Each iteration records duration per pipeline step \(\delta_{\omega,s}\)
(seconds).  Mean time per step:

\begin{equation}
  \bar{\delta}_s = \frac{1}{n_\mathrm{iter}} \sum_\omega \delta_{\omega,s}.
\end{equation}

Total expected time for $n_\mathrm{iter}$ iterations:

\begin{equation}
  T_\mathrm{total} \approx n_\mathrm{iter} \sum_s \bar{\delta}_s.
\end{equation}

\paragraph{Storage.}

Total storage is approximately

\begin{equation}
  D_\mathrm{total} \approx n_\mathrm{iter} \cdot D_\mathrm{run},
\end{equation}

where $D_\mathrm{run}$ is the size of one run directory (measured once, e.g.\
with \texttt{du -s run\_0001}).

% -----------------------------------------------------------------------
\section{Cohort Delta Protocol}

To avoid rebuilding centroids from scratch at every iteration,
MethylValidation can pass explicit cohort deltas to \texttt{methyl-centroid}:

\begin{equation}
  \text{effective\_samples}_\omega
  = (\text{samples}_{\omega-1} \setminus \text{remove}_\omega)
    \cup \text{add}_\omega,
\end{equation}

where $\text{samples}_{\omega-1}$ is the cohort of the previous iteration.
If no previous state is available (first iteration or unavailable centroid),
a full rebuild from the resolved cohort is performed.

% -----------------------------------------------------------------------
\section{Summary}

\begin{table}[h]
\centering
\begin{tabular}{@{}ll@{}}
\toprule
Aspect & Role \\
\midrule
Inputs & Sample lists, \texttt{train\_fraction}, \texttt{n\_iterations},
         base project, output base \\
Outputs & \texttt{all\_metrics.csv}, \texttt{metrics\_summary.json},
          \texttt{step\_timings.csv} \\
Metric distribution & Empirical: mean, std, percentiles per metric \\
Time estimation & Mean duration per step; total $= n_\mathrm{iter} \cdot \sum \bar\delta$ \\
Storage estimation & $n_\mathrm{iter} \times D_\mathrm{run}$ \\
\bottomrule
\end{tabular}
\caption{Summary of the MethylValidation Monte Carlo protocol.}
\end{table}

\end{document}
